% abstract_intro.tex — Abstract and Introduction
% NeurIPS 2026 submission: Spectral Conformal Prediction for Neural PDE Operators

\begin{abstract}
Neural operators such as Fourier Neural Operators (FNOs) learn PDE solution maps
efficiently but lack calibrated uncertainty quantification with finite-sample
guarantees. We introduce \emph{spectral nonconformity scores} for conformal
prediction on neural operators, defined as weighted sums over Fourier error
coefficients: $s_w(u) = \sum_k w_k |\hat{u}_{\mathrm{err}}(k)|^2$. This
formulation unifies $L^2$ scores (uniform weights, via Parseval's identity),
Sobolev scores (polynomial weights $(1+|k|^2)^s$), and learned
frequency-dependent weights under a single framework with distribution-free
coverage guarantees. Across 57 configurations (3 architectures $\times$ 5 seeds,
plus ablations over $\alpha$ and calibration set size) on real Darcy flow data,
learned spectral weights produce 42\% tighter prediction bands
than uniform scores on FNO while maintaining valid coverage across all
configurations. Non-spectral architectures exhibit $3.8\times$ wider bands,
confirming that frequency-structured scores exploit the spectral bias inherent
in Fourier-based operators. Our results establish spectral conformal prediction
as a principled, architecture-aware approach to uncertainty quantification for
neural PDE solvers.
\end{abstract}

\section{Introduction}
\label{sec:introduction}

Neural operators learn mappings between infinite-dimensional function spaces and
have achieved strong empirical performance on partial differential equations
(PDEs) including Navier--Stokes, Darcy flow, and Burgers' equation
\citep{li2021fourier,lu2021learning,kovachki2023neural}. Fourier Neural
Operators (FNOs) parameterize integral kernels in the frequency domain, enabling
resolution-invariant learning with quasi-linear complexity. Despite their
predictive accuracy, neural operators are deployed without calibrated uncertainty
estimates---a critical gap for safety-relevant applications in engineering
simulation, climate modeling, and medical imaging where practitioners must know
\emph{when} to trust a prediction.

Existing approaches to uncertainty quantification (UQ) for neural operators
inherit the limitations of standard deep learning UQ. Monte Carlo (MC) Dropout
\citep{gal2016dropout} requires dropout layers that FNO and Tensorized FNO
(TFNO) architectures typically lack, yielding degenerate estimates (0\% coverage
in our experiments). Deep ensembles \citep{lakshminarayanan2017simple} provide
useful heuristics but offer no finite-sample coverage guarantees and scale
poorly: training $M$ independent neural operators multiplies an already
substantial computational cost. Bayesian neural operator approaches
\citep{magnani2022approximate} introduce approximate inference but depend on
distributional assumptions that are difficult to verify for PDE solution
manifolds.

Conformal prediction \citep{vovk2005algorithmic,shafer2008tutorial} provides
distribution-free, finite-sample coverage guarantees under the sole assumption
of exchangeability. Given a calibration set and a \emph{nonconformity score}
measuring prediction quality, conformal prediction constructs prediction sets
that contain the true output with probability at least $1 - \alpha$ for any
user-specified $\alpha$. Recent work has applied conformal methods to regression
\citep{romano2019conformalized}, time series \citep{stankeviciute2021conformal},
and image segmentation \citep{angelopoulos2022image}. However, applying
conformal prediction to function-valued outputs requires nonconformity scores
defined over function spaces, and the standard choice---pointwise or global
$L^2$ residuals---treats all spatial frequencies uniformly, ignoring the
structure of neural operator errors.

Our key insight is that FNO errors are not uniformly distributed across
frequencies. The \emph{spectral bias} of neural networks
\citep{rahaman2019spectral} is amplified in Fourier-based operators: FNOs
process each frequency mode through learned weight matrices, producing errors
whose magnitude varies systematically across the spectrum. We exploit this
structure by defining spectral nonconformity scores
$s_w(u) = \sum_k w_k |\hat{u}_{\mathrm{err}}(k)|^2$, where the weights $w_k$
control sensitivity to different frequency bands. Uniform weights recover the
$L^2$ norm via Parseval's identity; polynomial weights
$w_k = (1 + |k|^2)^s$ yield Sobolev-type scores penalizing high-frequency
errors; and learned weights $w_k$ trained to minimize prediction band width on a
validation set produce the tightest bands while preserving the conformal coverage
guarantee. Because the weights affect only the scalar score and not the conformal
calibration mechanism, coverage validity is maintained for any fixed weight
choice.

\paragraph{Contributions.} We make the following contributions:
\begin{enumerate}[leftmargin=*,itemsep=2pt,topsep=2pt]
  \item \textbf{Spectral nonconformity scores.} We introduce a family of
    nonconformity scores defined via weighted Fourier coefficients of the
    prediction residual. We prove that conformal prediction with any fixed
    spectral weight vector $w$ satisfies the standard $1-\alpha$ marginal
    coverage guarantee (Theorem~\ref{thm:coverage}).

  \item \textbf{Learned spectral weights.} We show that optimizing the weight
    vector $w$ to minimize average prediction band width on a held-out set
    yields 42\% tighter bands than uniform weights on FNO (3.145 vs.\ 5.463
    mean band width) and 41\% tighter on TFNO (2.356 vs.\ 3.970), while
    maintaining valid coverage across all tested configurations.

  \item \textbf{Comprehensive empirical validation.} We evaluate across 57
    experiments on real Darcy flow data (64$\times$64, neuralop/Zenodo FEM
    solutions) spanning 3 architectures (FNO, TFNO, DeepONet), 4 significance
    levels, 4 calibration set sizes, and 5 random seeds. We confirm that
    non-spectral architectures (DeepONet) produce $3.8\times$ wider $L^2$ bands
    than FNO, providing direct evidence that spectral scores exploit the
    frequency-domain structure of Fourier-based operators.
\end{enumerate}
